\chapter{Conclusions}

\section{Summary of Work Done}
The group built a parcel tracking simulator in core Java. It tracks parcels
through four hubs, calculates the fee for three speeds, rejects invalid
tracking IDs, and saves its data in text files. Unit tests check these
features.

\section{Key Learnings and Outcomes}
The group learned the following from the project.

\textbf{Classes and encapsulation.} We learned to split a problem into small
classes: \texttt{Hub}, \texttt{Milestone}, \texttt{Route}, \texttt{Parcel}, and
others. Each class keeps its fields private and has one job. This made the code
easier to read and to change.

\textbf{Inheritance and abstract classes.} We learned to put what is common in
a base class and what is different in subclasses. The abstract method
\texttt{calculateTotalFee()} forces each speed class to give its own fee rule.

\textbf{Polymorphism.} We learned that one call can run different code. The
call to \texttt{calculateTotalFee()} in \texttt{Parcel} runs the code of the
real speed class, so \texttt{Parcel} does not need to know which speed it has.

\textbf{Factory and separation of work.} We learned to keep the choice of
class in one place, \texttt{ParcelFactory}. We also learned to keep the menu
(\texttt{Main}), the commands (\texttt{Tracker}), and the files
(\texttt{ManifestStore}) in separate classes.

\textbf{Exception handling.} We learned to write a custom exception, to throw
it where the error is found, and to catch it where the program can recover.
A wrong tracking ID now prints a message and the menu comes back.

\textbf{File input and output.} We learned to write a file again from the
start (the manifest) and to add to the end of a file (the status log). We also
learned to rebuild the objects from the files when the program starts, and to
close each file after use with try-with-resources, also when an error happens.

\textbf{Testing.} We learned to write unit tests that test one class at a
time, to check the expected value against the actual value, and to run a build
test before the unit tests. We learned to run all tests with one call.

The outcomes of the project are:
\begin{itemize}
\item A working command-line program of about 400 lines of core Java.
\item Three shipment speeds with their own delivery days and fees.
\item Data that stays after the program closes.
\item A build test and 50 unit tests with 109 checks that all pass.
\item This report.
\end{itemize}
